\chapter{Alcance y delimitaciones}

El proyecto abordó la reactivación progresiva del robot cuadrúpedo Nova Spot
Micro mediante modelado, simulación, control de locomoción, preparación
electrónica y diseño de una futura capa correctiva de aprendizaje por refuerzo.
El sistema se limitó a postura y locomoción de baja velocidad tipo paso y gateo
sobre una superficie plana y bajo condiciones controladas.

\section{Componentes desarrollados}

\begin{itemize}
\item Se formularon e implementaron modelos cinemáticos y dinámicos nominales,
jacobianos y descripciones URDF/SDF y MJCF.
\item Se integró una arquitectura ROS~2 con control articular, métricas,
supervisor, fase de marcha, contacto, IMU y margen de estabilidad nominal.
\item Se implementaron y evaluaron en simulación postura, marcha paso y gateo.
El galope se conservó únicamente como experimento opcional, deshabilitado por
defecto y excluido de los criterios principales.
\item Se preparó la Raspberry Pi, se verificó comunicación DDS e I$^2$C y se
realizaron pruebas progresivas del PCA9685 sin autorizar una marcha física.
\item Se definieron el protocolo, las métricas, las cantidades experimentales y
las semillas para la comparación final.
\end{itemize}

\section{Delimitaciones}

No se abordaron navegación autónoma, conectividad en la nube, percepción del
entorno ni locomoción rápida sobre terreno irregular. La política aprendida no
comandará directamente PWM ni sustituirá límites articulares, saturaciones o el
supervisor independiente.

El modelo actual se consideró nominal computable porque no se completó la
identificación física de geometría, masas, fricción, holgura y actuadores. El
margen de estabilidad calculado en simulación no se presentó como medición
física. Tampoco se reportó seguimiento articular real, ya que el PCA9685 no
proporciona realimentación de posición.

La seguridad eléctrica, calibración de doce servos, marcha física, entrenamiento
PPO y comparación nominal frente a nominal más RL permanecieron fuera de los
resultados cerrados de esta versión. Se mantuvieron como trabajo requerido para
cumplir integralmente los objetivos cuarto y quinto.

\section{Productos obtenidos y productos abiertos}

Se obtuvieron código versionado, modelos de simulación, documentación
matemática, controladores nominales, pruebas automatizadas, bolsas de datos,
analizadores, informes y un protocolo experimental. La entrega final deberá
añadir la caracterización física completa, diagramas y validación de seguridad,
calibraciones, políticas entrenadas, conjuntos de datos físicos y comparación
estadística final. La separación evita presentar productos esperados como si
hubieran sido obtenidos.
